% ============================================================
%  §IV  Implementation
%  Paper 3 — Constraint-Aware Ordering for Directed Traceability
%  ÉTS PhD / ACM Distributed Ledger Technologies, Autumn 2026
% ============================================================

\section{Implementation}
\label{sec:implementation}

We implemented the three architectures compared in the evaluation on top of
Hyperledger Fabric~v2.5, the de-facto reference platform for permissioned
blockchains.  The codebase consists of four independently deployable
components: a directed-traceability chaincode that encodes the resource
lifecycle, a use-case-agnostic constraint layer embedded in the Raft orderer
(\arch{M_D} endogenous), a lightweight HTTP coordinator paired with a standard
Fabric network (\arch{M_D} exogenous, ZooKeeper baseline), and a Caliper
workload driver that runs all three variants under identical network conditions.
The total implementation adds approximately 3\,200 lines of Go and
400 lines of JavaScript to the Fabric tree; no modifications to the Fabric peer
or the gossip/delivery layers are required.

%-------------------------------------------------------------
\subsection{Directed-Traceability Chaincode}
\label{subsec:impl-chaincode}

The chaincode (\texttt{integration/chaincode/directed}) implements the
resource lifecycle defined in Paper~2:
\textsc{RegisterResource}, \textsc{Transfer}, \textsc{Use}, \textsc{Publish},
and \textsc{Revoke}.  Each resource is stored under its primary key as a JSON
object carrying an identifier, type, current holder, subset-group membership,
usage conditions (including $K_\text{max}$), and an append-only trace log.

\paragraph{Constraint partitioning.}
The three usage constraints are enforced at different layers depending on their
scope:
\begin{itemize}
  \item \textbf{\(C_\text{auth}\)}: the proposing agent must be the current
    holder.  This is a \emph{local} constraint verifiable from the resource's
    own state; the chaincode enforces it by comparing the proposal creator
    identity against the stored holder at endorsement time.
  \item \textbf{\(C_\text{excl}\)}: at most one holder per resource at any
    instant.  Again local; the chaincode rejects any \textsc{Transfer} that
    would violate exclusivity.
  \item \textbf{\(C_\text{global}\)}: at most $K_\text{max}$ distinct agents
    may simultaneously hold resources within the same \emph{SubsetGroup}.
    This is a \emph{global} constraint: its evaluation requires a view of the
    concurrent holder set across \emph{all} resources in the group, which is
    distributed across ledger state entries and cannot be evaluated atomically
    by a single chaincode simulation.  It is therefore the sole constraint
    delegated to the orderer under \arch{M_D} endogenous.
\end{itemize}

This partitioning is the concrete realisation of the structural observation in
Paper~2, Theorem~2: $C_\text{global}$ is the constraint for which a
chaincode-only (\arch{M_L}) architecture produces a positive Intra-block
Violation Rate (IVR), because two conflicting transactions can be endorsed
simultaneously against an identical stale view.

\paragraph{World-state query function.}
To support the Init operation of Definition~1 (\S\ref{sec:model}), the
chaincode exposes a read-only function \textsc{GetAllResources}.  This function
performs a range scan over the primary-key space (lexicographic range
$[\texttt{0x20},\texttt{0x7E}]$, i.e.\ printable ASCII) and returns all
stored resources as a JSON array.  Composite-key entries, which Fabric stores
with a null-byte prefix, are excluded from this range and therefore not
returned.  The orderer uses this function to materialise the initial constraint
cache without replaying the block log (\S\ref{subsec:impl-snapshot}).

%-------------------------------------------------------------
\subsection{Generic Constraint Interface}
\label{subsec:impl-interface}

The orderer extension is designed to be \emph{use-case agnostic}: the
constraint layer knows nothing about biobank specimens, material-transfer
agreements, or any other domain.  Callers supply a concrete implementation of
the \texttt{Evaluator} interface, which encodes the constraint set $\mathcal{C}$
for their specific use case.

\paragraph{Core types.}
The \texttt{constraint} package (\texttt{orderer/consensus/etcdraft/constraint})
exposes three main abstractions:

\begin{itemize}
  \item \textbf{\texttt{CacheState}} — an opaque interface representing the
    constraint-relevant projection $\pi(S_\text{world})$ of the ledger world
    state.  It carries a single method \texttt{Clone()}, which produces a
    deep copy for use as the speculative state.

  \item \textbf{\texttt{TxView}} — the constraint-relevant projection of a
    single transaction envelope: its identifier \texttt{TxID}, the chaincode
    function name \texttt{Function}, its positional arguments \texttt{Args},
    and the channel identifier.  Parsing a raw Fabric envelope into a
    \texttt{TxView} is delegated to a \texttt{TxParser} supplied by each
    \texttt{Evaluator} implementation.

  \item \textbf{\texttt{Evaluator}} — the central interface, corresponding to
    Definition~1 of \S\ref{sec:model}:
    \begin{description}
      \item[\texttt{Init(snapshot)}] builds the initial stable cache state
        from a point-in-time world-state snapshot (§\ref{subsec:impl-snapshot}).
      \item[\texttt{Evaluate(s, tx)}] implements $\mathrm{Conf}(s, u_t)$:
        returns \emph{true} iff transaction \texttt{tx} is admissible in state
        \texttt{s}, i.e.\ all constraints in $\mathcal{C}$ are satisfied.
        It must not modify \texttt{s}.
      \item[\texttt{Apply(s, tx)}] computes $\delta(s, \mathrm{tx})$: the
        next cache state obtained by advancing \texttt{s} with an admitted
        transaction.  Called immediately after \texttt{Evaluate} returns
        \emph{true}.
      \item[\texttt{TxParser()}] returns the \texttt{TxParser} for this
        evaluator, used to extract a \texttt{TxView} from each raw envelope
        before \texttt{Evaluate} is called.
    \end{description}
\end{itemize}

\paragraph{WorldStateSnapshot interface.}
\begin{verbatim}
type WorldStateSnapshot interface {
    GetByRange(namespace, startKey, endKey string) ([]KeyValue, error)
}
\end{verbatim}
This interface abstracts the source of the initial world state consumed by
\texttt{Init}.  Decoupling the query mechanism from the evaluator enables
testing with in-process fakes while deploying with the live peer-endorsed
implementation described in \S\ref{subsec:impl-snapshot}.

%-------------------------------------------------------------
\subsection{Stable/Speculative Cache Management}
\label{subsec:impl-cache}

The \texttt{StateManager} struct implements the stable/speculative state pair
of Definition~2.  It is the only object in the orderer that calls
\texttt{Evaluator} methods; no other code in the chain touches the cache
directly.

\paragraph{Lifecycle.}
Four operations map directly onto the abstract operations of Definition~1:
\begin{enumerate}
  \item \textbf{Init} (called once at construction time): invokes
    \texttt{evaluator.Init(snapshot)} to build the initial $s_\text{stable}$
    and sets $s_\text{spec} = s_\text{stable}.\texttt{Clone()}$.
  \item \textbf{BeginBatch}: resets $s_\text{spec} \leftarrow
    s_\text{stable}.\texttt{Clone()}$, satisfying invariant~I2 of
    Definition~2 at the start of each block construction cycle.
  \item \textbf{ProcessBatch} (Algorithm~1): iterates over the envelope
    batch.  For each envelope the \texttt{TxParser} extracts a
    \texttt{TxView}; envelopes that do not yield a \texttt{TxView} (e.g.\
    configuration transactions) are admitted unconditionally.  For chaincode
    invocations, \texttt{Evaluate}$(s_\text{spec}, \mathrm{tx})$ is called.
    On \emph{true}, the envelope is admitted and $s_\text{spec}
    \leftarrow \texttt{Apply}(s_\text{spec}, \mathrm{tx})$ advances the
    speculative state.  On \emph{false}, the envelope is rejected and
    $s_\text{spec}$ is unchanged.
  \item \textbf{Commit}: promotes $s_\text{stable} \leftarrow s_\text{spec}$
    after the Raft quorum has confirmed the block.  A follower guard —
    \texttt{if !inBatch \{ return \}} — makes \texttt{Commit} a no-op on
    follower nodes that commit blocks received from the leader but never
    ran \texttt{BeginBatch}; this ensures invariant~I3 is not violated
    on the follower path.
  \item \textbf{Rollback}: restores $s_\text{spec} \leftarrow
    s_\text{stable}.\texttt{Clone()}$ on Raft-round failure, discarding the
    in-flight speculative state.
\end{enumerate}

\paragraph{Concurrency model.}
The block construction goroutine is the sole writer of the stable/speculative
pair under condition~C2 (Proposition~1): batch evaluation is strictly
sequential, one batch at a time.  The only concurrent caller is the Raft apply
goroutine, which calls \texttt{Commit} or \texttt{Rollback} after a consensus
round.  A single mutex guards this promotion path; no lock is held during
\texttt{Evaluate} or \texttt{Apply}.

%-------------------------------------------------------------
\subsection{World-State Snapshot via Peer Endorser}
\label{subsec:impl-snapshot}

The Init operation of Definition~1 requires materialising $\pi(S_\text{world})$
— the constraint-relevant projection of the current ledger state — at orderer
startup.  A naive approach would replay the block log from genesis, applying
the state-transition function $\delta$ to each committed transaction.  This
approach is proportional in cost to the \emph{log length} rather than the
\emph{state size}, and produces the correct result only for constraints that
depend on the net effect of the full history (rather than the current state).

We instead implement Init by querying a trusted peer's world state directly,
using Fabric's gRPC endorser protocol.  The orderer constructs a
\textsc{ProcessProposal} request invoking the directed chaincode's
\textsc{GetAllResources} function, sends it to a designated \emph{reader peer},
and parses the JSON array response.  The result represents the exact current
world state at the instant of the query, regardless of log length or the
presence of superseded intermediate states.  This design matches the description
in \S\ref{subsec:model-init}: Init takes a snapshot $\pi(S_\text{world})$, not
a block iterator.

\paragraph{PeerEndorserSnapshot.}
The concrete implementation lives in the \texttt{peersnapshot} package
(\texttt{orderer/consensus/etcdraft/constraint/peersnapshot}).  It implements
the \texttt{WorldStateSnapshot} interface through a single
\texttt{PeerEndorserSnapshot} struct that wraps a persistent gRPC connection
to the trusted peer:
\begin{verbatim}
func (p *PeerEndorserSnapshot) GetByRange(
        namespace, _, _ string) ([]KeyValue, error) {
    // Build ChaincodeInvocationSpec{GetAllResources}
    // Sign with reader identity (ECDSA, DER-encoded)
    // ProcessProposal → parse JSON array response
    // Return []KeyValue{Key: resource.ID, Value: raw JSON}
}
\end{verbatim}
The namespace parameter selects the chaincode to query; the start/end key
bounds are not forwarded to the peer, since \textsc{GetAllResources} returns
all primary-key entries unconditionally.

\paragraph{Reader identity.}
The orderer's own MSP identity (OrdererMSP) is not a member of application
channel policies; it cannot endorse chaincode proposals.  We therefore
configure a dedicated \emph{reader identity}: a regular application MSP
certificate and private key whose holder has read access to the channel.
The orderer loads this identity from the file system at startup (paths given by
environment variables \texttt{FABRIC\_CONSTRAINT\_READER\_CERT} and
\texttt{FABRIC\_CONSTRAINT\_READER\_KEY}) and uses it exclusively to sign the
\textsc{GetAllResources} proposal.  The reader identity carries no channel write
privileges and is used only at startup.  The proposal is signed with
SHA-256~+~ECDSA and the signature is DER-encoded in-process, without invoking
the Fabric MSP stack.

\paragraph{Configuration.}
The peer snapshot subsystem is configured via six environment variables on
the orderer process:
\begin{center}
\begin{tabular}{ll}
\toprule
Variable & Purpose \\
\midrule
\texttt{FABRIC\_CONSTRAINT\_PEER\_ADDR}   & Peer address (\texttt{host:port}) \\
\texttt{FABRIC\_CONSTRAINT\_PEER\_TLSCA}  & Path to PEM TLS CA certificate \\
\texttt{FABRIC\_CONSTRAINT\_READER\_CERT} & Path to PEM reader certificate \\
\texttt{FABRIC\_CONSTRAINT\_READER\_KEY}  & Path to PEM reader private key \\
\texttt{FABRIC\_CONSTRAINT\_READER\_MSP}  & Reader MSP identifier \\
\texttt{FABRIC\_CONSTRAINT\_CHAINCODE}    & Chaincode name (default: \texttt{directed}) \\
\bottomrule
\end{tabular}
\end{center}
Per-channel activation is controlled by the comma-separated variable
\texttt{FABRIC\_CONSTRAINT\_CHANNELS}: channels not listed in this variable
use the standard Fabric pipeline unchanged, preserving full backwards
compatibility.

%-------------------------------------------------------------
\subsection{Biobank Evaluator (\(C_\text{global}\) Instantiation)}
\label{subsec:impl-evaluator}

The \texttt{biobank} package
(\texttt{orderer/consensus/etcdraft/constraint/biobank}) is the concrete
instantiation of \texttt{Evaluator} for the directed-traceability use case.
It encodes only $C_\text{global}$; $C_\text{auth}$ and $C_\text{excl}$ are
local constraints already enforced by the chaincode at endorsement time and
need not be re-evaluated by the orderer.

\paragraph{Cache structure.}
The \texttt{BioankCacheState} maintains two maps:
\begin{itemize}
  \item \textbf{\texttt{resources}}: $\text{resourceID} \to (
    \text{currentHolder},\ \text{subsetGroup},\ K_\text{max},\ \text{status})$
    — the constraint-relevant projection of each resource.
  \item \textbf{\texttt{groups}}: $\text{subsetGroup} \to (\text{Holders},\
    K_\text{max})$ — the concurrent holder set and limit for each group.
    $|\text{Holders}|$ is the quantity directly compared against $K_\text{max}$
    during evaluation.
\end{itemize}
\texttt{Clone()} performs a full deep copy of both maps, including the
\texttt{Holders} set for each group, so that the speculative copy is
entirely independent of the stable copy.

\paragraph{Init.}
\texttt{BioankEvaluator.Init} calls \texttt{snapshot.GetByRange("directed","","")}
to obtain all resources, then projects each \emph{active} resource with a
non-empty \texttt{SubsetGroup} and a positive $K_\text{max}$ into the two maps.
Resources in terminal states (\texttt{inactive}) and resources without a group
or a concurrency limit are ignored, because they do not affect $C_\text{global}$.

\paragraph{Evaluate.}
Only \textsc{RegisterResource} and \textsc{Transfer} transactions affect the
concurrent holder count and therefore undergo $C_\text{global}$ evaluation.
All other functions (\textsc{Use}, \textsc{Publish}, \textsc{Revoke}) are
admitted unconditionally by the evaluator (they may still be rejected by the
chaincode).

For a \textsc{Transfer} of resource $r$ to agent $a$:
\[
  \mathrm{Conf}(s,\, \textsc{Transfer}(r, a))
  \;=\;
  \begin{cases}
    1 & \text{if } a \in \mathrm{Holders}(s,\, r.\text{group}) \\
    1 & \text{if } |\mathrm{Holders}(s,\, r.\text{group})| + 1 \leq K_\text{max} \\
    0 & \text{otherwise}
  \end{cases}
\]
If the target agent already holds another resource in the same group, the
holder count does not increase and the transfer is unconditionally admitted.
Otherwise, the count after admission must not exceed $K_\text{max}$.

For a \textsc{RegisterResource} of a new resource in group $g$ with limit
$K_\text{max}$ by agent $a$, the same logic applies with the registering agent
as the initial holder.

\paragraph{Apply.}
After a \textsc{Transfer} is admitted, \texttt{Apply} updates both maps:
(i)~it removes the previous holder from the group's \texttt{Holders} set
\emph{only if} they hold no other active resource in the same group, and
(ii)~it adds the new holder.  After a \textsc{Revoke} or \textsc{Publish}
the resource is marked inactive and its holder is similarly removed.

The \texttt{removeHolderFromGroup} helper scans all entries in
\texttt{resources} to determine whether the departing holder retains any
other resource in the same group before removing them from \texttt{Holders}.
This scan is bounded by the number of resources in the cache, which is at most
the total number of active resources in the channel — typically small enough to
run synchronously in the block construction goroutine.

\paragraph{TxParser.}
\texttt{BioankTxParser.Parse} unwraps the five-layer Fabric envelope hierarchy
(Envelope → Payload → Transaction → TransactionAction →
ChaincodeInvocationSpec) using Protocol Buffer unmarshalling to extract the
chaincode function name and its positional arguments as raw byte slices.
Configuration transactions (channel header type $\neq$ \textsc{EndorserTransaction})
return \texttt{nil}, causing the envelope to be admitted unconditionally.

%-------------------------------------------------------------
\subsection{Integration into the Raft Orderer}
\label{subsec:impl-raft}

The constraint layer is integrated into Fabric's \texttt{etcdraft.Chain} via
four targeted hook points.  No changes are made to the Raft consensus protocol
itself, the block storage layer, or the endorsement pipeline.

\paragraph{Chain extension.}
The \texttt{Chain} struct acquires a single new field:
\begin{verbatim}
constraintMgr constraintManager
\end{verbatim}
where \texttt{constraintManager} is a minimal interface exposing
\texttt{BeginBatch}, \texttt{ProcessBatch}, \texttt{Commit}, and
\texttt{Rollback}.  When this field is \texttt{nil} (the default), every hook
is a no-op and the chain behaves identically to a stock Fabric orderer.
The nil-check idiom preserves the standard pipeline with zero overhead for
channels that do not activate constraint-aware ordering.

\paragraph{Hook 1 — Block proposal (\texttt{propose}).}
The existing \texttt{propose} function, called by the block construction
goroutine immediately before proposing a new block to Raft, is extended with:
\begin{verbatim}
batch = c.filterBatch(batch)
if len(batch) == 0 { continue }
\end{verbatim}
\texttt{filterBatch} calls \texttt{BeginBatch} (cloning stable into speculative)
and then \texttt{ProcessBatch} (Algorithm~1), returning only the admitted
envelope subset.  If the entire batch is rejected, block proposal is skipped
for that cycle.  This hook implements condition~C2 of Proposition~1:
\texttt{ProcessBatch} is the unique call site of the evaluator within each
Raft round, ensuring sequential cache access.

\paragraph{Hook 2 — Block commit (\texttt{writeBlock}).}
Immediately after \texttt{support.WriteBlock} writes the confirmed block to
the ledger, the hook
\begin{verbatim}
if c.constraintMgr != nil { c.constraintMgr.Commit() }
\end{verbatim}
promotes $s_\text{spec}$ to $s_\text{stable}$.  Placement directly after
\texttt{WriteBlock} implements condition~C3 of Proposition~1: the stable cache
update is atomic with the ledger write within the Raft apply goroutine.

\paragraph{Hook 3 — Leader abdication (\texttt{becomeFollower}).}
When the Raft node steps down from leader to follower, any speculative state
built for in-flight proposals will never be committed.  The hook
\begin{verbatim}
if c.constraintMgr != nil { c.constraintMgr.Rollback() }
\end{verbatim}
discards this state by resetting $s_\text{spec} \leftarrow
s_\text{stable}.\texttt{Clone()}$, satisfying the Rollback operation of
Definition~1.

\paragraph{Hook 4 — Chain activation (\texttt{HandleChain}).}
The consenter's \texttt{HandleChain} method, called once per channel at orderer
startup, checks whether the channel appears in the \texttt{FABRIC\_CONSTRAINT\_CHANNELS}
environment variable.  If so, it executes the following initialisation
sequence:
\begin{enumerate}
  \item Read the peer snapshot configuration from environment variables
    (\texttt{peersnapshot.FromEnv}).
  \item Dial the trusted peer with mutual TLS and obtain a
    \texttt{PeerEndorserSnapshot} (\texttt{peersnapshot.New}).
  \item Construct a \texttt{BioankEvaluator} and pass it together with the
    snapshot to \texttt{constraint.NewStateManager}, which calls
    \texttt{Init} to build the initial stable cache.
  \item Wire the resulting \texttt{StateManager} into the chain via
    \texttt{chain.EnableConstraintAwareOrdering(mgr)}.
\end{enumerate}
If any step fails — unreachable peer, malformed TLS certificate, or Init
error — \texttt{HandleChain} returns an error and the orderer refuses to start
for that channel, preventing a channel from starting in an unconstrained state.

%-------------------------------------------------------------
\subsection{ZooKeeper Coordinator (M_D Exogenous Baseline)}
\label{subsec:impl-zk}

The exogenous baseline replicates the guarantee of $M_D$ without modifying the
orderer.  A stateless HTTP coordinator process (the ``ZooKeeper coordinator'',
\texttt{tools/zkcoordinator}) mediates all transfer requests: clients must
obtain a \emph{reservation} before submitting a transaction to Fabric, and
confirm or cancel it after the transaction commits or fails.

\paragraph{API.}
The coordinator exposes four endpoints over HTTP/1.1:
\begin{itemize}
  \item \textbf{POST /v1/reserve} — atomically checks $C_\text{global}$ against
    the in-memory counter and, if admission is possible, increments the count
    and returns a \texttt{reservationId}.
  \item \textbf{POST /v1/confirm} — removes a pending reservation and records
    the transfer as committed (decrements the outgoing holder's count if
    appropriate).
  \item \textbf{POST /v1/cancel} — releases a reservation without committing,
    decrementing the count previously incremented by \textsc{Reserve}.
  \item \textbf{POST /v1/release} — releases a holder slot explicitly
    (used after \textsc{Revoke} or \textsc{Publish}).
\end{itemize}
A background goroutine expires stale reservations (those not confirmed or
cancelled within a configurable timeout) every 10 seconds, preventing leaked
counts from client crashes.

\paragraph{Protocol and limitations.}
The client workload driver (\texttt{transfer-zk.js}) implements the full
three-phase protocol: Reserve → Fabric \textsc{Transfer} submit → Confirm
(or Cancel on error).  The coordinator enforces $C_\text{global}$ at the
Reserve step; IVR is zero as long as every client faithfully executes
Confirm/Cancel.  However, a client crash between Reserve and Confirm (or
Confirm and ledger commit) breaks this guarantee: the coordinator's counter
diverges from the ledger state.  The \texttt{expireStaleReservations}
mechanism mitigates — but does not eliminate — this divergence.  This
structural fragility, absent in the endogenous architecture, is one of
the empirical observations of the evaluation (\S\ref{sec:evaluation}).

\paragraph{Implementation.}
The coordinator is a single-binary Go service (no ZooKeeper dependency
is actually used; the name reflects the exogenous coordination role
described in Paper~2).  It maintains one in-memory \texttt{map[string]*GroupCounter}
protected by a mutex; persistence is intentionally omitted to match
the Paper~2 exogenous model, which assumes the coordinator restarts
cleanly or is HA-replicated externally.  The binary is packaged as a
multi-stage Docker image and deployed alongside the Fabric network
as a \texttt{docker-compose} sidecar.

%-------------------------------------------------------------
\subsection{Correctness of the Integration}
\label{subsec:impl-correctness}

We verify that the three conditions of Proposition~1 are satisfied by the
implementation:

\textbf{C1 — Leader uniqueness.}
Raft guarantees that at most one node holds the leader role in a given term.
The \texttt{constraintMgr} is activated via \texttt{EnableConstraintAwareOrdering}
on the leader's \texttt{Chain} instance; follower instances never call
\texttt{BeginBatch} or \texttt{ProcessBatch}.  The follower guard in
\texttt{Commit} (\texttt{if !inBatch \{ return \}}) ensures that blocks
received from the leader are written to the follower ledger without corrupting
the follower's stable cache, which was correctly initialised by \texttt{Init}
at startup.

\textbf{C2 — Sequential cache access.}
\texttt{filterBatch} (and therefore \texttt{BeginBatch} + \texttt{ProcessBatch})
is called exclusively from the block construction goroutine within
\texttt{propose}.  Fabric's etcdraft implementation processes at most one
outstanding proposal at a time (bounded by \texttt{MaxInflightBlocks});
no concurrent evaluation occurs.  The mutex in \texttt{StateManager} guards
only the \texttt{Commit}/\texttt{Rollback} path called from the Raft apply
goroutine.

\textbf{C3 — Atomic commit.}
\texttt{Commit} is called synchronously from \texttt{writeBlock}, immediately
after \texttt{support.WriteBlock} returns.  The Raft apply goroutine does not
yield between these two calls; no other thread can observe an intermediate
state in which the ledger has been updated but the stable cache has not.
\texttt{Rollback}, symmetrically, is called from \texttt{becomeFollower} before
the goroutine resumes normal follower operation.

Together these three conditions ensure that the sequential equivalence of
Proposition~1 holds: processing a batch of $n$ transactions through
\texttt{ProcessBatch} under the evolving speculative state is equivalent to
evaluating each transaction individually in the order it was admitted, against
the running state.  Theorem~1 then guarantees that the IVR of the endogenous
architecture is identically zero for any executions that satisfy C1--C3.

%-------------------------------------------------------------
\subsection{Lines of Code Summary}
\label{subsec:impl-loc}

\begin{table}[h]
\centering
\caption{Implementation effort by component.}
\label{tab:loc}
\begin{tabular}{lrr}
\toprule
Component & Go (lines) & JS (lines) \\
\midrule
Directed-traceability chaincode            & 620  & —   \\
Generic constraint interface (\texttt{constraint}) & 280  & —   \\
Biobank evaluator (\texttt{biobank})       & 455  & —   \\
Peer snapshot (\texttt{peersnapshot})      & 328  & —   \\
Orderer hook points (\texttt{chain}, \texttt{consenter}) & 95 & — \\
ZooKeeper coordinator                      & 490  & —   \\
Unit tests (constraint, chain)             & 415  & —   \\
Caliper workload drivers                   & —    & 410 \\
\midrule
\textbf{Total}                             & \textbf{2\,683} & \textbf{410} \\
\bottomrule
\end{tabular}
\end{table}

The orderer hook points account for fewer than 100 lines, reflecting the
minimal invasiveness of the integration: the entire constraint enforcement
mechanism is expressed as four nil-checked hooks on the existing \texttt{Chain}
struct, with no changes to the Raft protocol, block storage, or client-facing
gRPC services.
